\documentclass[10pt]{amsart}

\usepackage[margin=1in]{geometry}
\usepackage[T1]{fontenc}
\usepackage{lmodern}
\usepackage{amsmath,amssymb,amsthm,mathtools}
\usepackage{microtype}
\usepackage{booktabs}
\usepackage{url}
\usepackage{array}
\usepackage{enumitem}
\usepackage[dvipsnames]{xcolor}
\usepackage[colorlinks=true,linkcolor=MidnightBlue,citecolor=OliveGreen,urlcolor=MidnightBlue]{hyperref}
\usepackage[capitalize,nameinlink]{cleveref}

\setlist{leftmargin=*,itemsep=2pt,topsep=4pt}
\emergencystretch=2.5em
\raggedbottom
\setlength{\tabcolsep}{4pt}
\renewcommand{\arraystretch}{1.12}

\theoremstyle{definition}
\newtheorem{definition}{Definition}[section]
\theoremstyle{remark}
\newtheorem{remark}[definition]{Remark}
\numberwithin{equation}{section}

\newcommand{\A}[1][]{\textup{[A\if\relax\detokenize{#1}\relax\else, #1\fi]}}
\newcommand{\Oa}[1][]{\textup{[O\if\relax\detokenize{#1}\relax\else, #1\fi]}}
\newcommand{\HC}{\#\mathrm{HC}}
\newcommand{\tst}{t^{*}}
\newcommand{\nst}{n^{*}}
\newcommand{\Hh}{\textup{[H]}}
\DeclareUrlCommand\lf{\urlstyle{tt}}
\newcommand{\sz}{\operatorname{size}}

\begin{document}

\title[When do the new \#SAT algorithms beat brute force?]{When do the new \#SAT algorithms beat brute force?\\
Crossover points of the proven bounds\\ behind the ETH refutations}

\author{Vaibhav Tulsyan}

\date{November 2026}

\subjclass[2020]{68Q25, 68W40, 68Q17}
\keywords{Exponential Time Hypothesis, model counting, fine-grained complexity, explicit constants, galactic algorithms}

\begin{abstract}
In October 2026 OpenAI and Anthropic claimed deterministic algorithms that count the satisfying
assignments of a $k$-CNF in time $2^{O(n\log\log n/\sqrt{\log n})}$ or similar, refuting the
Exponential Time Hypothesis. Neither manuscript computes from which $n$ on its \emph{proven} bound is
smaller than brute force. We compute this crossover $\nst$ for the bounds as they are proved,
keeping every explicit term and constant and treating existential constants as parameters, and
report $\tst=\log_2\nst$. For Anthropic's official Lean main theorem on $3$-CNFs with $m=4.26n$
clauses, $\tst=552\,040$ if the unknown polynomial degree $q$ of its Turing-machine time bound is
$1$, and $\tst$ grows roughly like $q^2$; re-optimising the free parameters of the same lemmas gives
$245\,009$. For OpenAI's interpolation algorithm, $\tst\approx3.02\cdot10^{13}$ as written and
$\approx8.53\cdot10^{7}$ with re-optimised parameters, mostly by using an exact binomial tail in
place of a Hoeffding window. Our exact algorithm for counting Hamiltonian cycles,
which rests on Anthropic's manuscript, beats Held--Karp on cubic graphs from $\tst=55\,440$.
Anthropic's $3$-SAT decision machine searches exhaustively for every $n<2^{2^{1722}}$. In every
Anthropic-based case $\nst$ is a staircase, fixed by the first recursion level at which a single
clique exponent falls below $n$. These numbers are computations about proven upper bounds, not
theorems and not running times; they are conditional on the claimed refutations, which are under
review, and nothing practical follows.
\end{abstract}

\maketitle

% =====================================================================
\section{Introduction}\label{sec:intro}
% =====================================================================

The Exponential Time Hypothesis (ETH) of Impagliazzo and Paturi~\cite{ImpagliazzoPaturi01}
asserts that $3$-SAT on $n$ variables needs time $2^{\delta n}$ for some $\delta>0$. In October
2026 two manuscripts, by OpenAI~\cite{OpenAI26} and by Anthropic~\cite{Anthropic26}, each claimed a
deterministic algorithm that counts the models of a $k$-CNF in subexponential time, with Lean~4
formalisations of the top-level statements. We write \A{} and \Oa{} for them. The Main Theorem of
\A{} gives time $2^{O(\min\{km/\sqrt{\log m},\;k^2n\log\log n/\sqrt{\log n}\})}$; \Oa{} gives
$(n+m)^{O_k(1)}2^{O_k(n\log\log n/\sqrt{\log n})}$ bit operations, and $2^{O_{k,D}(n/\sqrt{\log
n})}$ when $m\le Dn$. These are \emph{claimed} results, under review. In a companion
paper~\cite{Tulsyan26} we derived from \A{} an exact deterministic algorithm counting the
Hamiltonian cycles ($\HC$) of a graph with $n$ vertices and $m=O(n)$ edges in time
$2^{O(n/\sqrt{\log n})}$; the Lean development and the scripts used here are in the
repository~\cite{TulsyanRepo26}.

An $O(\cdot)$ in an exponent says nothing about when the algorithm first wins. The question is
not idle: the sources fix almost all of their constants (the Lean development of \A{} fixes every
schedule parameter), so the bound they prove is a concrete function of $n$, up to a few
existential constants. Neither manuscript computes where its concrete bound crosses brute force (the
Lean development of \A{} records only that its decision machine searches exhaustively below
$2^{2^{518}}$, see \cref{sec:corrections}), and we are not aware of such a computation; we do it here.

\begin{definition}\label{def:cross}
Let $T_{\mathrm{alg}}(n)$ be a proven upper bound on the number of steps of an algorithm on inputs
with $n$ variables (or vertices), with every explicit term and polynomial factor kept and the free
parameters chosen as the tier says (\cref{sec:tiers}), and let $T_{\mathrm{bf}}(n)$ be the
brute-force bound: $2^n$ for $\#3$-SAT and $n^22^n$ (Bellman, Held--Karp~\cite{Bellman62,HeldKarp62})
for $\HC$. The \emph{crossover} $\nst$ is the least $n$ with $T_{\mathrm{alg}}(n)<T_{\mathrm{bf}}(n)$.
Since $\nst$ has between $10^3$ and $10^{14}$ binary digits (and $2^{1747}$ for one machine), we report $\tst=\log_2\nst$.
\end{definition}

We write $\rho_{\mathrm{alg}}(n)=\log_2T_{\mathrm{alg}}(n)/n$ for the \emph{rate}; for $\#3$-SAT
the crossover is where the rate first drops below $1$. \Cref{tab:headline} is the summary.

\begin{table}[ht]
\centering\footnotesize
\caption{Crossover $\tst=\log_2\nst$ of the proven bounds. A: as written; B: same lemmas, free
parameters re-optimised; C: constant changes justified outside the formal proofs; \Hh: heuristic.
Rows 1--3 use $q=1$ for the unknown degree of the Turing-machine time bound (the explicit-part
lower bound, \cref{sec:explicit}); for other $q$ see \cref{tab:q}. In rows 7--8 the values in
parentheses use the degree network of~\cite{TulsyanRepo26}.}\label{tab:headline}
\begin{tabular}{@{}lllll@{}}
\toprule
& algorithm (brute force) & A & B & C \\
\midrule
1 & \A, official Lean thm., $\#3$-SAT, $m=4.26n$ ($2^n$) & $552\,040$ & $245\,009$ & $119\,103$; $31\,358$ \Hh\\
2 & \A, $n$-route, $\#3$-SAT, any $m$ ($2^n$) & $3.42\cdot10^{10}$ & $1.61\cdot10^{10}$ & $3.93\cdot10^{9}$; $1.05\cdot10^{9}$ \Hh\\
3 & \A, $3$-SAT decision machine ($2^n$) & $2^{1747}$ & --- & ---\\
4 & \Oa, Engine I, $\#3$-SAT, any $m$ ($2^n$) & $3.02\cdot10^{13}$ & $8.53\cdot10^{7}$ & $3.82\cdot10^{7}$\\
5 & \Oa, Engine II ($2^n$) & \multicolumn{3}{l}{only a lower bound: $\log_2(\text{wires})\ge2^{85}$ ($2^{510}$ at full depth)}\\
6 & $\HC$, cubic graphs ($n^22^n$) & $55\,440$ & $44\,356$ & $20\,888$; $4\,905$ \Hh\\
7 & $\HC$, $\Delta=4$ ($n^22^n$) & $112\,860$ & $90\,291$ & $43\,205$; $10\,470$ ($9\,799$) \Hh\\
8 & $\HC$, $\Delta=8$ ($n^22^n$) & $458\,100$ ($112\,860$) & $366\,482$ ($90\,292$) & $177\,700$ ($49\,161$);\\
&&&& $46\,182$ ($19\,350$) \Hh\\
\bottomrule
\end{tabular}
\end{table}

So, for example, the official Lean theorem of \A{} first certifies a win over $2^n$ on $3$-CNFs at
the satisfiability threshold density at about $n=2^{552\,040}$, and only if $q=1$; with $q=2$ it is
$2^{2.21\cdot10^6}$ and with $q=10$ it is $2^{3.54\cdot10^7}$. OpenAI's Engine I as written needs
$n\approx2^{3.02\cdot10^{13}}$. Our $\HC$ algorithm, which has no Turing-machine degree because it
counts word operations, has the smallest crossover of all.

\subsection*{What the numbers are}
Each $\tst$ is the output of an exact computation (\cref{sec:repro}) applied to formulas that we
transcribed from the sources (file and lemma references below). The transcription is ours, the
computations are not formally verified, and none of the numbers is a theorem, in Lean or on paper.
Where a number depends on a constant that the sources only assert to exist, we say so and report its
sensitivity. A crossover of a proven bound says where the \emph{guarantee} first beats brute force;
the algorithms might do better, and the guarantees might be improvable, which is what tiers B and
C measure.

% =====================================================================
\section{Method}\label{sec:method}
% =====================================================================

\subsection{Exact evaluation in the log domain}\label{sec:logdomain}
All quantities are represented by their base-$2$ logarithms, evaluated with
mpmath~\cite{mpmath} at $60$ significant digits. The level sizes $K_\ell=2^{L_\ell}$ have
$L_\ell$ up to about $10^{9}$ and are never built as integers; the exponents $L_\ell$, the level
rules and the Lean functions \lf{Nat.size}, \lf{Nat.clog}, floors and ceilings are evaluated
exactly on these representations. We put $n=2^t$ with real $t$, and $m=\alpha n$ for the
$m$-route. Within one band of $t$ on which the recursion level is fixed the rate decreases in $t$,
and the crossover is found by bisection in $t$ to relative precision $10^{-6}$; tier~B and~C
crossovers are minimised over all levels. Terms reported as ``below $10^{-x}\cdot n$'' are
printed by the script in decimal notation.

\subsection{The schedule of \texorpdfstring{\A}{[A]}}\label{sec:schedule}
The algorithm of \A{} is a recursion over levels $\ell$. Its Lean development fixes the schedule
(\lf{Levels.Gen}, \lf{statements/LevelsStatement.lean}; all Lean paths are relative to the
\lf{proof/} directory of \A):
\begin{equation}\label{eq:schedule}
K_0=2^{x_0},\quad g_\ell=2^{3\ell+x_0}K_\ell,\quad K_{\ell+1}=g_\ell^2,\quad \zeta_\ell=2^{-\ell-5},\quad
e_0=\tfrac12,\quad e_{\ell+1}=(\rho+3\zeta_\ell)\,e_\ell+\frac{3}{K_\ell},
\end{equation}
with $\rho=1/\sqrt2$. The official route uses $x_0=18$, the least value allowed
(\lf{maintheorem/CountedRoute/Defs.lean}); then $\log_2K_\ell=60\cdot2^\ell-6\ell-42$, which the
script asserts for $\ell<40$ together with condition~(ii) of the grouping lemma,
$K_\ell\ln(2g_\ell K_\ell)\le2\zeta_\ell^2g_\ell$. The number $e_\ell$ is the clique exponent: the
clique algorithm at level $\ell$ runs in $2^{e_\ell K^2\lambda}(K(\lambda+B))^{C_0K^2}$ steps
(\lf{statements/CliqueLevelsStatement.lean}). The official main theorem, however, does not use
$e_\ell$ but the rational number $e'_\ell$ returned by \lf{num}/\lf{den} in
\lf{rates/GeneralPieces/Num/Between.lean} (lines~34--37): $e'_\ell=2^{-\ell/2}$ for even $\ell$
and $1.4\cdot2^{-(\ell+1)/2}$ for odd $\ell$ (\cref{sec:corrections}).

\begin{table}[ht]
\centering\small
\caption{The Lean schedule at $x_0=18$ (the last column is the ratio of the printed values).}\label{tab:schedule}
\begin{tabular}{@{}rrllll@{}}
\toprule
$\ell$ & $\log_2K_\ell$ & $e_\ell$ & $e'_\ell$ & $e_\ell/\rho^\ell$ & $e'_\ell/e_\ell$\\
\midrule
0 & 18 & 0.5 & 1.0 & 0.5 & 2.0\\
2 & 186 & 0.3019 & 0.5 & 0.6038 & 1.66\\
4 & 894 & 0.1586 & 0.25 & 0.6342 & 1.58\\
6 & 3\,762 & 0.08026 & 0.125 & 0.6421 & 1.56\\
8 & 15\,270 & 0.04026 & 0.0625 & 0.6441 & 1.55\\
9 & 30\,624 & 0.02848 & 0.04375 & 0.6444 & 1.54\\
10 & 61\,338 & 0.02014 & 0.03125 & 0.6446 & 1.55\\
12 & 245\,646 & 0.01007 & 0.01563 & 0.6447 & 1.55\\
\bottomrule
\end{tabular}
\end{table}

The ratio $e_\ell/\rho^\ell$ tends to $0.644772=\frac12\prod_j(1+3\zeta_j/\rho)$: the terms
$3\zeta_\ell$ cost a factor $1.29$ over $\rho^\ell/2$, while the terms $3/K_\ell$ are below
$2^{-17}$ at $\ell=0$ and doubly exponentially smaller afterwards. Two facts drive everything
below: $e_\ell$ halves every two levels, and $\log_2K_\ell$ doubles with every level.

\subsection{Tiers and labels}\label{sec:tiers}
\emph{Tier~A} is the bound as written: the Lean definitions for \A, the paper's parameter choices
for \Oa{} and for~\cite{Tulsyan26}. \emph{Tier~B} uses the same proven lemmas with their free
parameters re-optimised for each $n$ (for \A: the level and the numerator $e_\ell$ in place of
$e'_\ell$, both allowed by other Lean routes of \A; for \Oa: the parameters $s,t,r,Q,R,\tau$
subject only to the numerical hypotheses of \lf{lem:layout}). \emph{Tier~C} changes constants in
ways justified outside the formal proofs. C1 re-tunes $\zeta_\ell=2^{-(\ell+z)}$ and takes the least
$g_\ell/K_\ell$ that condition~(ii) allows; this rests on the written proof of \A, not on Lean. For
\Oa, tier C replaces the bound $m_{\mathrm{test}}\le6n$ by $m_{\mathrm{test}}\le n+n/\log_2(4/3)+2\approx3.41n$, which the proofs of \lf{lem:sparse-checks} and
\lf{lem:greedy-hashes} give but the paper does not state. Items marked \Hh{} are heuristic: C2 and
C3 assume a base case of the clique recursion with exponent $e_0=5/18$ ($=2.5/9$) or
$e_0=\omega/9=0.2635$ with $\omega=2.371552$~\cite{VXXZ24}, as a fast-matrix-multiplication base
case might give; C4 adds $x_0=10$ (which needs the pieces constant $c\le1017$); C5 drops the
$K^g$ grouping search. None of the \Hh{} changes is proved, in the sources or here.

\subsection{Explicit and existential constants}\label{sec:constants}
\Cref{tab:constants} lists every constant that is existential in Lean or hidden in an $O(\cdot)$ of
the sources, and how it is handled. No existential constant is ever given a value; each is a
parameter, and we report either a table in it or the \emph{break-even} value it would have to
exceed to move $\nst$.

\begin{table}[ht]
\centering\small
\caption{Existential and hidden constants.}\label{tab:constants}
\begin{tabular}{@{}>{\raggedright\arraybackslash}p{3.0cm}>{\raggedright\arraybackslash}p{5.6cm}>{\raggedright\arraybackslash}p{6.6cm}@{}}
\toprule
constant & where & effect on $\nst$\\
\midrule
$q$, degree of the TM time bound \A & \lf{Computes} (\lf{variables/PairsToBound/Defs.lean}, l.~192): steps $\le q(|x|+2^{\mathrm{expo}}+1)^q$; folded into the $C$ of the final forms & \textbf{the only one that matters}: crossover needs rate $<1/q$; $\tst$ grows about $4\times$ per doubling of $q$ (\cref{tab:q})\\
$C_0$, clique factor \A & $\forall c\,\exists C_0$ in \lf{ExplicitCost.RecursionSolves}; enters \lf{needOf} as $C_0K^2\sz(\cdot)$ & none: tables for $C_0\in\{0,1,10\}$ coincide; at the tier-A $\nst$ of row~1 it matters only if $C_0>2^{429\,342}$\\
$C_1$, grouping factor (Lemma~8 of \A) & $\exists C_1$ in \lf{GroupingExplicit}; absent from the Lean pad exponent & none for \A; for $\HC$ (cubic, tier A) only if $C_1>2^{55\,421.7}$; $C_0$ there only if $C_0>2^{51\,725.6}$\\
$C_2$, the $O(\log n)$ in $\log\tau$ ($\HC$) & \cite{TulsyanRepo26}, \lf{research/exact-hc-theorem.md} (6.1)--(6.2) & none: only if $C_2>2^{55\,421.0}$\\
$n_0\le n_1$, ``large $n$'' of the $n$-route \A & \lf{ThmVariablesAForEveryMachine}: $\exists n_0\,\forall n_1\ge n_0$ & cannot be bounded; can only increase $\nst$, so the $n$-route values are lower bounds\\
$(n+m)^{O_k(1)}$ and other polynomial factors \Oa & throughout & none: a factor $(n+m)^c$ matters only if $c$ is of order $2^{\lambda^*}$, $\lambda^*\approx3.0\cdot10^{13}$\\
Engine II constants $O_C(\cdot)$, $O_{C,A}(\cdot)$ \Oa & \lf{thm:tensor-contraction} & not boundable; only an explicit parameter threshold is reported\\
\bottomrule
\end{tabular}
\end{table}

\subsection{Explicit-part lower bounds}\label{sec:explicit}
Setting $C_0=C_1=C_2=0$ and $q=1$ and ignoring $n_0,n_1$ only decreases the bounds, so the
resulting $\tst$ is a rigorous lower bound on the crossover of the proven bound; we also drop
$\log_2q$ and the input length $|x|$ in $q(|x|+2^{\mathrm{expo}}+1)^q\ge2^{q\cdot\mathrm{expo}}$,
which again errs downwards. For the padded Lean machines of \A, which write a pad of
$2^{\mathrm{expo}}$ units and fall back to exhaustive search when $\mathrm{expo}>n$ (\lf{planM},
\lf{plan}), it is also a lower bound on where the machine itself leaves exhaustive search. For
\Oa{} Engine~I and for $\HC$ the remaining constants are so far from their break-even values that
the lower bound \emph{is} the crossover of the proven bound; for the SAT routes of \A{} it is the
value at $q=1$, and the true value is the $q$-dependent one of \cref{tab:q,tab:nroute}.

% =====================================================================
\section{Results}\label{sec:results}
% =====================================================================

\subsection{\texorpdfstring{\A}{[A]}, the official main theorem (\texorpdfstring{$m$}{m}-route)}\label{sec:mroute}
\lf{Solution.lean} (line~44) defines \lf{MainResults.mainTheorem} as
\lf{CountedRoute.mainTheoremHolds}. Its exponent in $m$ is \lf{needTop}
(\lf{rates/ClauseBound/Alg/Top.lean}): with $K=K_\ell$ for the largest $\ell$ with $K_\ell^9\le m$
(\lf{root9}), $g=3K^8$, $P=\max(1,2k\lceil m/g\rceil)$ and $\lambda=\lfloor(K^3+2kmK+gP)/K^3\rfloor$,
\begin{equation}\label{eq:needTop}
\mathrm{needTop}=\max\Bigl(3K^8\sz K,\;P\bigl(\tfrac{2g}{K}+1\bigr),\;(K-1)\lambda,\;
\lceil e'_\ell K^2\lambda\rceil+C_0K^2\sz\bigl(K(\lambda+B+1)\bigr)+2\Bigr),
\end{equation}
the last term being \lf{needOf} of \lf{rates/CountedToPadded/Defs.lean} (line~148) with
$B=2gP/K$. The machine runs in at most $q(|x|+2^{\mathrm{needTop}}+1)^q$ steps, and searches
exhaustively when $\mathrm{needTop}>n$. \Cref{tab:mrate} shows the rate at $q=1$; \cref{tab:mcross}
the crossovers; \cref{tab:q} the dependence on $q$.

\begin{table}[ht]
\centering\small
\caption{Rate $\mathrm{needTop}/n$ of the official theorem, tier A, $m=4.26n$, $q=1$, $C_0=0$, at
$n=2^{2^j}$. From $j=8$ on, the solver term \texttt{needOf} is the maximum in~\eqref{eq:needTop}.}\label{tab:mrate}
\begin{tabular}{@{}rrrrrrrr@{}}
\toprule
$j$ & 8 & 12 & 16 & 18 & 20 & 24 & 30\\
level $\ell$ & 0 & 3 & 6 & 8 & 10 & 14 & 20\\
rate & 25.56 & 8.946 & 3.195 & 1.597 & 0.7987 & 0.1997 & 0.02496\\
\bottomrule
\end{tabular}
\end{table}

\begin{table}[ht]
\centering\small
\caption{Crossover $\tst$ of the $m$-route, $k=3$, $m=\alpha n$, $q=1$; $\ell^*$ is the level at the
crossover for $\alpha=4.26$. C1 is justified by the written proof of \A, not by Lean.}\label{tab:mcross}
\begin{tabular}{@{}lrrrr@{}}
\toprule
tier & $\alpha=1$ & $\alpha=3$ & $\alpha=4.26$ & $\ell^*$\\
\midrule
A: \lf{CountedRoute}, \lf{root9}, $e'_\ell$ ($C_0\in\{0,1,10\}$) & $33\,858$ & $275\,614$ & $552\,040$ & 10\\
B: same lemmas, $e_\ell$, level free & $7\,163$ & $122\,176$ & $245\,009$ & 9\\
C1: $\zeta_\ell$ and $g_\ell/K_\ell$ re-tuned & $6\,459$ & $59\,247$ & $119\,103$ & 8\\
C2: C1 $+$ $e_0=5/18$ \Hh & $1\,449$ & $15\,405$ & $31\,358$ & 6\\
C3: C1 $+$ $e_0=\omega/9$ \Hh & $1\,353$ & $14\,412$ & $29\,341$ & 6\\
C4: C2 $+$ $x_0=10$ \Hh & $1\,161$ & $13\,036$ & $26\,621$ & 6\\
C5: C2 without the $K^g$ grouping search \Hh & $1\,265$ & $13\,473$ & $27\,431$ & 6\\
\bottomrule
\end{tabular}
\end{table}

\begin{table}[ht]
\centering\small
\caption{Sensitivity of $\tst$ to the Turing-machine degree $q$, $m$-route, $\alpha=4.26$.}\label{tab:q}
\begin{tabular}{@{}lrrrrr@{}}
\toprule
tier & $q=1$ & $q=2$ & $q=3$ & $q=5$ & $q=10$\\
\midrule
A & $5.52\cdot10^5$ & $2.21\cdot10^6$ & $4.42\cdot10^6$ & $8.85\cdot10^6$ & $3.54\cdot10^7$\\
B & $2.45\cdot10^5$ & $9.82\cdot10^5$ & $1.97\cdot10^6$ & $3.93\cdot10^6$ & $1.57\cdot10^7$\\
C1 & $1.19\cdot10^5$ & $4.78\cdot10^5$ & $9.58\cdot10^5$ & $3.55\cdot10^6$ & $1.42\cdot10^7$\\
C2 \Hh & $3.14\cdot10^4$ & $1.27\cdot10^5$ & $2.72\cdot10^5$ & $9.58\cdot10^5$ & $3.83\cdot10^6$\\
\bottomrule
\end{tabular}
\end{table}

\emph{Dominant term.} At every tier the cost at $\nst$ is the solver term
$\lceil e'K^2\lambda\rceil\approx2e'_\ell km$, that is, the clique exponent times the network
total $U\approx2km$. At the tier-A crossover ($\alpha=4.26$, $\ell=10$) it equals $0.79875\,n$,
while the grouping search, $P(2g/K+1)$ and $(K-1)\lambda$ are each below $10^{-18459}\cdot n$. So
the rate is $2e'_\ell k\alpha$ to that accuracy, and $\nst$ is fixed by two discrete facts: the first
level with $q\cdot2e'_\ell k\alpha<1$ (here $e'_{10}=1/32$, giving $25.56/32=0.79875$), and the
least $n$ at which that level may be used. In tier A this is the rule $K_\ell^9\le m$, so
$\tst=9\log_2K_{10}-\log_2 4.26=552\,042-2.09$, the value in \cref{tab:mcross}; in tier B the
binding constraint is the search term $3K^8\sz K\le n$, so $\tst\approx8\log_2K_9$. Since
$\log_2K_\ell$ doubles per level and $e_\ell$ halves every two levels, $\nst$ is a staircase: a
change of constant helps only once it crosses a level boundary, and then it halves $\tst$. Each
doubling of $q$ costs two levels, hence the factor $4$ in \cref{tab:q} (the steps at $q=3$ and
$q=5$ are the intermediate odd levels).

\subsection{\texorpdfstring{\A}{[A]}, the \texorpdfstring{$n$}{n}-route (any number of clauses)}\label{sec:nroute}
For formulas with arbitrarily many clauses, \A{} first sparsifies and then runs a network whose
total is $U\le6k(n+L_\Delta)+28kn\lceil\log_2(96\Delta L_\Delta)\rceil+2(3K^8)^2$, with
$L_\Delta=\lceil\log_2 2\Delta\rceil$ and $\Delta=\min(\text{\lf{occBound}},n+1)$
(\lf{variables/VariableBound/Net/Defs.lean}, \lf{variables/RecursionToBound/Defs.lean}). Here
\lf{occBound} $=2(\theta_0+\theta_1+\theta_2)+2=32\alpha^2+4\alpha+6$ with
$\alpha=\lceil8\varepsilon^{-1}\lceil\log_2(128/\varepsilon)\rceil\rceil$
(\lf{sparsification/SparsificationLemma/Cal/Defs.lean}) and the sparsification parameter
$\varepsilon=e'_\ell$ (\lf{epsOf}). The exponent is the solver exponent for this network plus a walk
term $\lceil\varepsilon n\rceil+2+\sz(\text{stack})$, at the level given by $K_\ell^{17}\le n$
(\lf{root17}).

\begin{table}[ht]
\centering\small
\caption{$n$-route: rate at $n=2^{2^j}$ (tier A, $q=1$), and crossovers $\tst$.}\label{tab:nroute}
\begin{tabular}{@{}lrrrrr@{}}
\toprule
$j$ & 22 & 28 & 34 & 37 & 40\\
level $\ell$ & 12 & 18 & 24 & 27 & 30\\
rate & 55.42 & 8.076 & 1.153 & 0.4251 & 0.1621\\
$U/n$ & 3\,546 & 4\,134 & 4\,722 & 4\,974 & 5\,310\\
\midrule
tier & $q=1$ & $q=2$ & $q=5$ & $q=10$ & \\
\midrule
A (\lf{root17}, $\varepsilon=e'_\ell$), $\ell^*=25$ at $q=1$ & $3.42\cdot10^{10}$ & $1.37\cdot10^{11}$ & $1.10\cdot10^{12}$ & $4.38\cdot10^{12}$\\
B ($e_\ell$, level free) & $1.61\cdot10^{10}$ & $6.44\cdot10^{10}$ & $5.15\cdot10^{11}$ & $2.06\cdot10^{12}$\\
C1 (re-tuned, $\varepsilon$ decoupled from $e$) & $3.93\cdot10^{9}$ & $1.68\cdot10^{10}$ & $1.26\cdot10^{11}$ & $5.37\cdot10^{11}$\\
C2 (C1 $+$ $e_0=5/18$) \Hh & $1.05\cdot10^{9}$ & $7.28\cdot10^{9}$ & $3.36\cdot10^{10}$ & $2.33\cdot10^{11}$\\
\bottomrule
\end{tabular}
\end{table}

At the tier-A crossover ($q=1$, $\ell=25$) the cap is $\Delta=2.8049\cdot10^{13}$ and $U/n=4806$;
the solver term is $0.82134\,n$, the walk term $1.709\cdot10^{-4}\,n$, and every $K$-dependent
term is below $10^{-6\cdot10^8}\cdot n$. The dominant term is again $e'_\ell U$, and
$\tst=17\log_2K_{25}$. Lean ties the sparsification parameter to the clique exponent; decoupling
them (tier C1, best at $\varepsilon\approx128e_\ell$) is worth about two levels, because the walk
term $\varepsilon n$ is cheap while $U$ depends on $\varepsilon$ through $\log_2\Delta\approx
2\log_2(1/\varepsilon)$. The existential threshold $n_1$ can only increase these values.

\subsection{\texorpdfstring{\A}{[A]}, the \texorpdfstring{$3$}{3}-SAT decision machine}\label{sec:decision}
The rate machine \lf{Sat3.inRateTM} (\lf{rates/PaddedDescent/}) chooses
$s=\lfloor3(\lfloor\log_2\lfloor\log_2n\rfloor\rfloor-911)/19\rfloor$ (\lf{sOf},
\lf{Params.lean} line~56) and searches exhaustively while $s<\text{\lf{s0K}}\,3=128$, i.e.\ for every
$n<2^{2^{1722}}$. Beyond that it writes a pad of $2^{2^{131}(\lfloor n/2^s\rfloor+1)}$ units
($\text{\lf{C0K}}\,3=8\cdot2^{128}=2^{131}$), and its time $q\cdot(\text{pad})^q$ is below $2^n$ only once
$q\cdot2^{131-s}<1$:
\begin{center}\small
\begin{tabular}{@{}lrrrr@{}}
\toprule
$q$ & 1 & 2 & 10 & 1000\\
needs $s\ge$ & 132 & 133 & 136 & 142\\
$\nst$ & $2^{2^{1747}}$ & $2^{2^{1754}}$ & $2^{2^{1773}}$ & $2^{2^{1811}}$\\
\bottomrule
\end{tabular}
\end{center}
This is a decision procedure, not a counter, and it is far worse than the counting routes because
of its padding constant.

\subsection{Our \texorpdfstring{$\HC$}{\#HC} algorithm}\label{sec:hc}
The $\HC$ algorithm of~\cite{Tulsyan26} is conditional on \A{} (Theorem~A(b), Lemma~8 and the
existence of its pieces). It counts word-RAM operations, so there is no Turing-machine degree and
no padding. We use the explicit bounds recorded in the repository~\cite{TulsyanRepo26}
(\lf{research/exact-hc-theorem.md}, (6.1)--(6.2)): $gP\le2.33(2m+n)+3g^3$, $U\le4.65m+g^3$,
$B\le2(2m+n)/g+5$, $\log\tau\le1.5(2m+n)/g+3+C_2\log n$, at most $n^{4m/\Theta}$ branches with
$\Theta=\max(2,\lceil n/g\rceil)$, and the level rule $K_\ell^{30}\le n$. The log of the bound is
\begin{multline*}
1+\max\bigl\{g\log K,\;e_\ell(U+gP/K+K^2)+C_0K^2\log(K(L+B'))+2gP/K+\log\tau\bigr\}\\
+C_1\log(g(P+B))+C_2\log n+\log\#\text{branches},
\end{multline*}
compared with $\log_2(n^22^n)$. The degree network (\lf{research/hc-degree-theorem.md}) has
$U\le2n\log_2(2e\bar d)+g^3$ and no branching; for $\Delta$-regular graphs we take $\bar d=\Delta$.

\begin{table}[ht]
\centering\small
\caption{$\HC$ versus $n^22^n$: crossover $\tst$, $m=\Delta n/2$. Tier A uses $K_\ell^{30}\le n$, the
Lean schedule and $e_\ell$; $C_0=C_1=C_2\in\{0,10\}$ give identical values. The network
of~\cite{Tulsyan26} has $U\approx4.65m$, the degree network $U\approx2n\log_2(2e\Delta)$.}\label{tab:hc}
\begin{tabular}{@{}lrrrrrr@{}}
\toprule
& \multicolumn{3}{c}{network of~\cite{Tulsyan26}} & \multicolumn{3}{c}{degree network}\\
\cmidrule(lr){2-4}\cmidrule(l){5-7}
tier & $\Delta=3$ & $\Delta=4$ & $\Delta=8$ & $\Delta=3$ & $\Delta=4$ & $\Delta=8$\\
\midrule
A & $55\,440$ & $112\,860$ & $458\,100$ & $55\,440$ & $112\,860$ & $112\,860$\\
B & $44\,356$ & $90\,291$ & $366\,482$ & $44\,357$ & $90\,291$ & $90\,292$\\
C1 & $20\,888$ & $43\,205$ & $177\,700$ & $40\,037$ & $43\,205$ & $49\,161$\\
C2 \Hh & $4\,905$ & $10\,470$ & $46\,182$ & $9\,080$ & $9\,799$ & $19\,350$\\
C3 \Hh & $4\,615$ & $9\,799$ & $43\,206$ & $9\,078$ & $9\,797$ & $19\,349$\\
C4 \Hh & $4\,330$ & $8\,695$ & $39\,079$ & $8\,021$ & $8\,024$ & $15\,846$\\
\bottomrule
\end{tabular}
\end{table}

\emph{Dominant term.} At the tier-A crossover for $\Delta=3$ ($\ell=5$) the term $e_\ell U$ is
$0.78846\,n$; $(e+2)gP/K$ is below $10^{-555}\cdot n$, $\log\tau$ below $10^{-4450}\cdot n$, the
grouping search and the branches below $10^{-12232}\cdot n$. The values of tier A are exactly
$30\log_2K_\ell$ ($30\cdot1848$, $30\cdot3762$, $30\cdot15\,270$). In tier B the binding lower-order
term is the copy-bit term $g^3=27K^{24}$ in $U$ and $gP$, so $\tst\approx24\log_2K_\ell$. The degree
network has the smaller $U$ from $\Delta\ge4$ ($8.885n$ against $9.30n$ at $\Delta=4$; at $\Delta=3$,
$8.06n$ against $6.98n$), but because levels are discrete it lowers tier~A only at $\Delta=8$, by a
factor $4.06$.

\subsection{\texorpdfstring{\Oa}{[O]}, Engine I (the interpolation algorithm)}\label{sec:engineI}
We call Engine I the interpolation algorithm of \lf{prop:interpolation-bound} of \Oa{} (line
numbers refer to \lf{main_full_version.tex}), which counts in
$(n+m)^{O_k(1)}2^{O_k(n(\log\log n)^2/\sqrt{\log n})}$ bit operations. Its proof sparsifies with
$\sigma=1/s$, branches on variables occurring in more than $b_0=\lceil M/n\rceil t$ scopes, and runs
the common counting routine (\lf{prop:general-count}, l.~2013), whose cost is a sum of three terms
(\lf{eq:three-costs}, l.~2055). With $n=2^\lambda$ the rate is
\[
\underbrace{\tfrac1s}_{\text{leaves}}+\underbrace{\tfrac1t}_{\text{branches}}
+\underbrace{\tfrac{\log_2C_q}{n}}_{\text{Taylor}}+\underbrace{2\tau\log_2(2D+1)}_{\text{measurements}}
+\underbrace{\tfrac2r}_{\text{group}}+\frac{\mathrm{polylog}}{n},
\qquad C_q=\sum_{s'\le q}2^{s'}\binom n{s'},\quad q=\Bigl\lfloor\frac{d_H}{Q}\Bigr\rfloor,
\]
with $d_H=kR(m_{\mathrm{test}}+2n)=24Rn$ for $k=3$, $m_{\mathrm{test}}\le6n$. We bound
$\log_2C_q\le q+nH_2(q/n)+\log_2(q+1)$ for $q<2n/3$ and $C_q\le3^n$ otherwise. Tier A takes the
parameters of \lf{lem:incidence-parameters} (l.~2747): $\lambda=\lceil\log_2n\rceil$,
$\nu=\lceil\log_2\lambda\rceil$, $r=\lfloor\sqrt\lambda/\nu^2\rfloor$,
$Q=\lfloor\lambda/(128\nu)\rfloor$, $R=(\beta+3)\lceil\sqrt\lambda\rceil$ with $\beta$ least with
$b,w\le\lambda^\beta$, $\tau=1/(r\lambda)$, $D=r^{2Q}$, and the sparsification constants of
\lf{app:sparsify}; if any numerical check of the proof fails ($r\ge2$, $Q>R$, $R\ge32$, $6R\le Q$,
$2(b+w)D^2\le M'$, $2Qr^2\le n'$, the tail bound), the algorithm searches exhaustively.

\begin{table}[ht]
\centering\small
\caption{Engine I, tier A. For $\lambda\le2^{36}$ a check fails ($s<2$ up to $2^{18}$; then $r<2$ or
$Q\le R$; then $6R>Q$) and the algorithm searches exhaustively.}\label{tab:oaiA}
\begin{tabular}{@{}rrrrrr@{}}
\toprule
$\log_2\lambda$ & 38 & 42 & 44 & 46 & 50\\
$\beta$ & 7 & 6 & 6 & 6 & 6\\
$r$ & 363 & 1\,188 & 2\,166 & 3\,964 & 13\,421\\
$q/n$ & 2.23 & 0.554 & 0.290 & 0.152 & 0.0412\\
rate & 1.593 & 1.548 & 1.16 & 0.7662 & 0.2892\\
\bottomrule
\end{tabular}
\end{table}

All checks first pass at $\lambda\approx8.07\cdot10^{10}$, where the rate is $1.599$. The rate first
drops below $1$ at $\lambda^*=\tst\approx3.0166\cdot10^{13}$ ($\log_2\lambda^*=44.778$), with
$\beta=6$, $r=2712$, $Q=5.24\cdot10^9$, $R=4.94\cdot10^7$, $q/n=0.2265$. The terms there are Taylor
$0.9984$, group $7.4\cdot10^{-4}$, leaves $3.7\cdot10^{-4}$, measurements $2.9\cdot10^{-6}$ and
branches $1.8\cdot10^{-7}$; the polylogarithmic terms (Tanner search, $3n(2m)^{P_h}$ calls) are
below $10^{-9\cdot10^{12}}$. \emph{The dominant term is the Taylor (Hasse-derivative) term
$\log_2C_q$}: the order $q=24Rn/Q$ is large because the window $R=(\beta+3)\lceil\sqrt\lambda\rceil$
comes from a Hoeffding bound~\cite{Hoeffding63} with a union over $b\le\lambda^\beta$ incidences,
which is far wider than the binomial tail requires.

\emph{Tier B} keeps the lemmas and only the numerical hypotheses of \lf{lem:layout}, and optimises
$s$, $t$, $r$, $R$, with $Q$ at its largest allowed value and $\tau=r\,b\Pr[W\notin I]$ the least
allowed, $W\sim\mathrm{Bin}(Q,(r-1)/(2r))$. The tail is bounded by the Chernoff--Kullback--Leibler
bound, which is a rigorous upper bound on the exact binomial tail. The search runs in floating
point (coordinate descent with restarts), and each reported point is re-evaluated in mpmath by the
same cost function as tier A.

\begin{table}[ht]
\centering\small
\caption{Engine I, tier B: best rate found at $\lambda=2^j$, and its parameters.}\label{tab:oaiB}
\begin{tabular}{@{}rrrrrr@{}}
\toprule
$j$ & 24 & 26 & 28 & 32 & 40\\
rate & 1.5083 & 1.0699 & 0.7097 & 0.2807 & 0.0349\\
$s$ & 128 & 128 & 128 & 512 & 4\,096\\
$t$ & 512 & 512 & 512 & 2\,048 & 16\,384\\
$r$ & 18 & 18 & 29 & 86 & 946\\
$R$ & 17\,150 & 34\,948 & 66\,152 & 244\,286 & 3\,463\,792\\
\bottomrule
\end{tabular}
\end{table}

The tier-B crossover is $\lambda^*_B\approx8.53\cdot10^7$ ($\log_2\lambda^*=26.35$), at which the
mpmath rate is $0.99983$ (Taylor $0.888$, group $0.100$, leaves $0.0078$, branches $0.0020$,
measurements $0.0018$): \emph{the exact-tail window shrinks $\tst$ by a factor of about
$3.5\cdot10^5$}. In tier B the binding trade-off is between the group term $2/r$ and the Taylor
term, through $Q_{\max}\propto\lambda/\log r$. Tier C ($m_{\mathrm{test}}\approx3.41n$, so
$d_H\approx16.2Rn$) gives $\lambda^*_C\approx3.82\cdot10^7$ ($\log_2\lambda^*=25.19$, rate $0.99985$).
We did not use two further savings that the proofs do not give: halving the group term, and a lower
precision than $\ell=n+1$.

\subsection{\texorpdfstring{\Oa}{[O]}, Engine II (tensor contraction)}\label{sec:engineII}
The sharper bounds of \Oa{} come from \lf{thm:tensor-contraction} (\lf{quantitative-reduction.tex},
l.~26), all of whose rate constants are $O_C(\cdot)$. What is explicit is the parameter regime of
\lf{quantitative-contraction.tex}: $\varepsilon=2^{-17}$, $u_d=\lfloor\varepsilon\lambda/2^d\rfloor$,
and the requirement $\lambda^{4/5}\le u_d$ at every active call, where now $\lambda$ is
$\log_2$ of the number of wires. Active calls reach depth $d<(\log_2C+\frac12\log_2\lambda)/3$,
and there the requirement reads $\varepsilon\lambda^{5/6}/C^{1/3}\ge\lambda^{4/5}$, i.e.\
$\lambda\ge2^{510}C^{10}$; at depth $0$ alone it needs $\lambda\ge\varepsilon^{-5}=2^{85}$. Reading the
parameter bounds as hypotheses that the proof needs, Engine II as written says nothing below
$\log_2(\text{wires})\approx2^{85}$, which is far beyond Engine I's crossover, and it gives no
finite upper bound on $\nst$ at all. Engine I is therefore the operative crossover of \Oa.

\subsection{Which levers matter}\label{sec:levers}
For \Oa{} one lever dominates: the exact binomial window in place of Hoeffding (tier B), a factor
$3.5\cdot10^5$ in $\tst$; the polynomial factors cannot matter. For \A{} the levers are, in order:
\begin{enumerate}[label=\textup{(\roman*)}]
\item \emph{the degree $q$}, existential and unbounded in Lean: about $\times4$ in $\tst$ per doubling;
\item \emph{a fast base case} ($e_0=5/18$, \Hh): $e_\ell$ times $0.556$, about $1.7$ levels,
$119\,103\to31\,358$ (factor $3.8$) in row~1;
\item \emph{$e_\ell$ in place of $e'_\ell$ and a free level} (tier B, proven in other Lean routes):
$552\,040\to245\,009$ (factor $2.25$);
\item \emph{re-tuning $\zeta_\ell$ and $g_\ell/K_\ell$} (C1): $245\,009\to119\,103$ (factor $2.06$);
Lean's $\zeta_\ell=2^{-(\ell+5)}$ has nearly the right shape, but $g_\ell=2^{3\ell+x_0}K_\ell$ is
larger than condition~(ii) needs;
\item minor: $x_0=10$ (\Hh, factor $1.18$) and dropping the grouping search (\Hh, factor $1.14$); $3/K_\ell$, $C_0$,
$C_1$, $C_2$ are irrelevant.
\end{enumerate}
From A to C2 the gain in row~1 is a factor $17.6$. For $\HC$ the same schedule levers apply.

% =====================================================================
\section{Four facts about the sources}\label{sec:corrections}
% =====================================================================
A first reading of the statements suggests four things that the sources, read in full, do not
say. We record what they do say, because each changes a number above.

\begin{enumerate}[label=\textup{(\arabic*)}]
\item \emph{The decision machine's floor is $2^{2^{1722}}$, not $2^{2^{518}}$.} The docstring of
\lf{statements/ThreeSatRateStatement.lean} (line~64) says that the machine of the proof is
exhaustive search for every $n<2^{2^{518}}$. This is true but weak: with \lf{sOf}, \lf{s0K} and
\lf{C0K} of \lf{rates/PaddedDescent/Params.lean} (lines~56, 83, 87) the machine searches
exhaustively for every $n<2^{2^{1722}}$, and its pad exceeds $2^n$ until $n\ge2^{2^{1747}}$, as
\lf{rates/PaddedDescent/Main.lean} itself states (lines~52--59).
\item \emph{The official theorem uses $e'_\ell\approx1.55\,e_\ell$, not $e_\ell$.} The official
\lf{MainResults.mainTheorem} (\lf{Solution.lean}, line~44) is \lf{CountedRoute.mainTheoremHolds},
whose solver \lf{solverOfMachine} pads to \lf{needOf}$=\lceil e'K^2\lambda\rceil+C_0K^2\sz(K(\lambda+B+1))+2$
(\lf{rates/CountedToPadded/Defs.lean}, line~148) with $e'=e'_\ell$ from \lf{num}/\lf{den}
(\lf{rates/GeneralPieces/Num/Between.lean}, lines~34--37). Since $e'_\ell\approx\rho^\ell$, it is
about $1.55\,e_\ell$ from $\ell=6$ on (\cref{tab:schedule}). The interface \lf{NumberBetween}
(\lf{statements/CliqueLevelsStatement.lean}) allows any number in $[e_\ell,2e_\ell]$; tier~B uses
$e_\ell$, its lower end. Separately, the same file says that in Turing-machine steps the exponent
$e_\ell K_\ell^2\lambda$ is claimed only up to a constant factor fixed with the machine (docstring,
lines~25--33); in counted operations (\lf{SolvesInFactor}) it is exact. We count the former through
$q$.
\item \emph{OpenAI's occurrence bound is explicit.} \lf{lem:sparsify} of \Oa{} (l.~3543) states
$B=O_k([\sigma^{-1}\log(2/\sigma)]^{k-1})$, but adds that $B$ is ``an explicitly computed integer'',
and \lf{app:sparsify} (``The size of the tree'', l.~3808) computes it: $\Gamma=2(3^{k-1}-1)$, $M$
the least power of $2$ with $M\ge\Gamma$ and $\Gamma\lceil\log_2(4CM)\rceil\le\sigma M$,
$C=4\sum_{j\le k-2}(3M)^j$, $B=2\sum_{j<k}M^j$. We use these values; no hidden constant enters
through $B$.
\item \emph{Anthropic's occurrence bound is explicit.} The $n$-route caps the number of
occurrences per variable by \lf{occBound} (\lf{variables/RecursionToBound/Defs.lean}, line~115),
$2\sum_j\theta_j+2=32\alpha^2+4\alpha+6$ for $k=3$, with $\alpha$, $\beta_c$, $\theta_c$ defined in
\lf{sparsification/SparsificationLemma/Cal/Defs.lean}. It is an explicit function of $\varepsilon$
and involves no existential constant.
\end{enumerate}

% =====================================================================
\section{Limitations}\label{sec:limitations}
% =====================================================================

\begin{enumerate}[label=\textup{(\alph*)}]
\item \emph{Everything is conditional.} The bounds evaluated here are those of the claimed
refutations \A{} and \Oa, and, for $\HC$, of our paper~\cite{Tulsyan26}, which assumes three
internal results of \A. Expert review of the manuscripts is pending. If their proofs fail, the
numbers lose their meaning.
\item \emph{Computations, not theorems.} The formulas were transcribed by us from Lean definitions
and from the \TeX{} sources (for \Oa, only the \lf{.tex} files were read) and evaluated by a Python
script. Neither the transcription nor the script is formally verified, and no number in this note
is machine-checked. The Lean development of \Oa{} certifies no finite crossover at all.
\item \emph{Existential constants.} The degree $q$ of \A's Turing-machine bound is unknown and
unbounded in Lean; the values of rows~1--3 of \cref{tab:headline} are those at $q=1$, i.e.\ lower
bounds. The thresholds $n_0\le n_1$ of the $n$-route can only increase its crossover. The
constants $C_0,C_1,C_2$ and the polynomial degrees of \Oa{} are shown to be irrelevant only within
the stated break-even ranges.
\item \emph{Heuristic tiers.} The entries marked \Hh{} (fast base case, $x_0=10$, no grouping
search) rest on changes that nobody has proved. Tier C1 rests on the written proof of \A{} and
tier C of \Oa{} on our reading of two proofs; neither is in Lean.
\item \emph{Engine II} of \Oa{} yields only a lower bound; its rate constants are $O_C(\cdot)$,
and the threshold $2^{85}$ uses our reading of its parameter bounds as hypotheses of the proof.
\item \emph{Proven bounds, not running times.} A crossover of a proven upper bound says nothing
about the actual running time, which may be far smaller; tier B is only the best point found by a
coordinate search, not a proven optimum. The $\HC$ cost model uses the explicit bounds of our
research notes, which accompany but are not identical with the Lean statements of~\cite{Tulsyan26}.
\item \emph{Nothing practical follows.} The smallest crossover in this note, $\tst\approx4\,330$
for $\HC$ with all heuristic changes, is an input with more than $10^{1300}$ vertices.
\end{enumerate}

% =====================================================================
\section{Reproducibility}\label{sec:repro}
% =====================================================================
The computation is \lf{research/scripts/crossover.py} in the repository~\cite{TulsyanRepo26},
with its notes in \lf{research/crossover.md}. It is run as \texttt{python3 -I crossover.py}
(isolated mode, so that no file of the working directory is imported) and writes
\lf{crossover.out} beside itself; it needs only the Python standard library and
mpmath~\cite{mpmath} at $60$ significant digits (\texttt{mp.mp.dps = 60}). On 3~November 2026, with
Python~3.13.17 and mpmath~1.3.0, a fresh run took $29$~seconds of wall-clock time and reproduced
the stored \lf{crossover.out} byte for byte. Every number in this note is either printed in that
file or, where stated, a ratio or difference of printed values. The script also asserts the
closed form of $\log_2K_\ell$ and condition~(ii) for the Lean schedule for $\ell<40$.

\section*{Acknowledgements}
This work was done with AI assistance (Claude Code, Anthropic): the transcription of the bounds,
the script and the text were produced in sessions with Claude Code. The numbers have been checked
against a fresh run of the script and the formulas against the sources as described, but the note
awaits human expert review. Responsibility for any errors rests with the author.

\bibliographystyle{amsplain}
\bibliography{refs}

\end{document}
